\section{Discussão e ameaças à validade}
\label{sec:discussao}
\subsection{Validade de construto}
Uma release pública é uma aproximação de implantação, e falhas de workflows
não identificam necessariamente incidentes em produção. Runs podem incluir
testes instáveis, falhas de infraestrutura ou de configuração. Recuperação
observada em CI não equivale à restauração de serviço percebida pelo usuário.
As variantes operacionais e as classes fixadas pelo enunciado devem ser comparadas somente
após integração e validação das métricas correspondentes.

\subsection{Validade interna e confiabilidade}
Limites de paginação, indisponibilidade de repositórios e retenção ou remoção
de históricos podem produzir dados incompletos mesmo numa janela fechada.
Releases podem ser editadas ou publicadas posteriormente; uma janela fechada
não garante imutabilidade dos registros. Cache preserva o snapshot observado,
mas deve ser distinguido de uma consulta atual. O default branch é observado
na coleta e pode ter sido diferente no passado. A censura e recuperações
anteriores ao início da janela exigem cuidado ao interpretar as durações.

\subsection{Validade externa}
Repositórios populares, usuários de Actions e projetos com releases frequentes
não representam todos os projetos de software. Selecionar os primeiros elegíveis
na ordem de estrelas introduz viés para projetos populares. As conclusions devem
ser limitadas à população e à janela estudadas. A discussão empírica das hipóteses
será adicionada após a auditoria, sem alterar seu registro anterior à coleta.

\subsection{Validade de conclusão}
A análise futura deve tratar comparações múltiplas, efeitos de subgrupos pequenos
e durações censuradas. As hipóteses complementares foram registradas depois do
piloto, o que limita sua interpretação como previsões anteriores aos dados.
Nenhum teste estatístico ou veredito das hipóteses é apresentado nesta sprint.
